\documentclass[12pt]{article}
\usepackage[utf8]{inputenc}
\usepackage[portuguese]{babel}
\usepackage{graphicx}
\usepackage{float}
\usepackage{geometry}
\usepackage{caption}
\usepackage{amsmath}
\captionsetup{labelfont=bf}
\geometry{a4paper, margin=2.5cm}

\title{Identificação de Dígitos através de Características Extraídas de Sinais de Áudio}
\author{Simão Tomás Botas Carvalho \\ nº2021223055}
\date{}

\begin{document}

\begin{figure}[H]
    \centering
    \includegraphics[width=0.8\textwidth]{uc_capa.png}
\end{figure}

\begingroup
\let\newpage\relax
\maketitle
\endgroup

\newpage


\section{Introdução}

Este relatório corresponde à Meta 2 do projeto de Análise e Transformação de Dados, que tem como
objetivo a identificação automatizada de dígitos falados (de 0 a 9, em inglês) a partir de gravações
de voz do dataset AudioMNIST.

Nesta fase, o trabalho foca-se na análise tempo-frequência dos sinais de áudio, através da
Transformada de Fourier de Curto Prazo (STFT) e da Transformada de Wavelet Discreta (DWT),
bem como na extração de características correspondentes e na aplicação de um classificador
automático.

A Meta 1 forneceu a base de trabalho: a estrutura de dados com os sinais pré-processados e as
características temporais e espectrais já calculadas. Nesta meta, essa estrutura é estendida com
novas características tempo-frequência, culminando na classificação dos dígitos e na avaliação do
desempenho do sistema.

\section{Metodologia Geral}

Tal como na Meta 1, o trabalho foi desenvolvido em MATLAB. A estrutura de dados (tabela) criada
anteriormente foi carregada a partir do ficheiro \texttt{.mat} guardado no final da meta anterior,
e foi sendo atualizada ao longo dos exercícios com as novas características extraídas.

O script principal da Meta 2 é o \texttt{main2\_main.m}. Foi também desenvolvido o script auxiliar
\texttt{optimize\_stft\_params.m}, que permitiu identificar os melhores parâmetros STFT para o
trabalho.

\section{Meta 2 -- Análise Tempo-Frequência e Classificação}

\subsection{Carregamento da Estrutura de Dados}

O script começa por carregar o ficheiro \texttt{meta1\_fourier\_and\_spectral\_features.mat},
gerado no final da Meta 1, que contém os sinais pré-processados e todas as características
calculadas anteriormente (temporais e espectrais). Isto permite retomar o trabalho sem necessidade
de reprocessar os dados originais.

\subsection{Transformada de Fourier de Curto Prazo (STFT)}

Para cada sinal de áudio, foi calculada a STFT (Short-Time Fourier Transform), que permite analisar
como o conteúdo espectral do sinal evolui ao longo do tempo. A STFT é definida como:

\[
X(m, \omega) = \sum_{n=-\infty}^{+\infty} x[n]\, w[n - m]\, e^{-j\omega n}
\]

onde $x[n]$ é o sinal de entrada, $w[n]$ é a janela de análise e $m$ é o índice temporal do
segmento. O resultado é um espectrograma de potência $P(f, t) = |X(m, \omega)|^2$.

Para calcular a STFT foi utilizada a função \texttt{spectrogram} do MATLAB. Foram testadas
diferentes parametrizações -- tamanhos de janela, sobreposições e número de pontos de FFT --
através do script \texttt{optimize\_stft\_params.m}, que avaliou sistematicamente combinações de:

\begin{itemize}
    \item Tamanho de janela: 256, 512, 1024, 2048 amostras;
    \item Sobreposição: 25\%, 50\%, 75\%;
    \item Número de pontos FFT (NFFT): 256, 512, 1024, 2048.
\end{itemize}

Os parâmetros que produziram melhor desempenho na classificação foram:

\begin{itemize}
    \item Janela de Hann com tamanho $w = 256$ amostras;
    \item Sobreposição de 50\% ($\text{noverlap} = 128$);
    \item $\text{NFFT} = 256$ pontos.
\end{itemize}

A figura seguinte mostra os espectrogramas obtidos para um exemplo de cada dígito (0 a 9) com
estes parâmetros.

% -----------------------------------------------------------------------
% FIGURA 1: Espectrogramas dos dígitos 0 a 9 (subplot 5x2)
% Introduzir aqui a figura gerada pelo bloco %% 15 do main2_main.m
% Nome sugerido para o ficheiro: stft_spectrograms.png
% -----------------------------------------------------------------------
\begin{figure}[H]
    \centering
    \includegraphics[width=\textwidth]{stft_spectrograms.png}
    \caption{Espectrogramas (STFT) para os dígitos de 0 a 9, com janela de 256 amostras,
             sobreposição de 50\% e NFFT=256.}
    \label{fig:stft}
\end{figure}

Observando os espectrogramas, é possível verificar diferenças no padrão temporal e espectral entre
dígitos. Dígitos com sons mais curtos e percussivos (como ``one'' ou ``eight'') apresentam
concentrações de energia mais localizadas no tempo, enquanto dígitos mais longos (como ``seven''
ou ``zero'') apresentam estruturas espectrais mais estendidas. A componente de baixa frequência
é geralmente mais energética, o que é consistente com a natureza dos sons vogais e das fricativas
presentes nos dígitos.

\subsection{Extração de Características Tempo-Frequência}

A partir da STFT, foram extraídas características globais do espectrograma e características
localizadas em três regiões temporais (início, meio e fim do sinal).

\subsubsection{Divisão Temporal}

Para capturar a evolução temporal do sinal, cada espectrograma é dividido em três regiões de tempo
de igual duração. Seja $N_T$ o número total de frames temporais do espectrograma e
$F_s = \lfloor N_T / 3 \rfloor$ o tamanho de cada região. As regiões são:

\begin{itemize}
    \item Região 1 (\textit{start}): frames $1$ a $F_s$;
    \item Região 2 (\textit{mid}): frames $F_s+1$ a $2F_s$;
    \item Região 3 (\textit{end}): frames $2F_s+1$ a $N_T$.
\end{itemize}

Para cada região $r$, considera-se a submatriz de potência $P_r$ e o espectro médio da região:
\[
\bar{P}_r[k] = \frac{1}{|r|} \sum_{t \in r} P[k, t].
\]

\subsubsection{Características Globais}

Para cada sinal, são calculadas as seguintes características globais com base no espectro médio
$\bar{P}[k] = \frac{1}{N_T}\sum_{t=1}^{N_T} P[k,t]$ e na potência total
$P_{\text{tot}} = \sum_k \bar{P}[k]$:

\begin{itemize}
    \item \textbf{Energia espectral} (\textit{energySpec}):
    \[
    E_{\text{spec}} = \sum_{k,t} P[k,t]
    \]

    \item \textbf{Frequência de pico} (\textit{peakFreq}):
    \[
    f_{\text{peak}} = f_{\,\arg\max_k \bar{P}[k]}
    \]

    \item \textbf{Centroide espectral} (\textit{specCentroid}):
    \[
    C = \frac{\sum_k f_k\, \bar{P}[k]}{P_{\text{tot}}}
    \]

    \item \textbf{Largura de banda espectral} (\textit{specBW}):
    \[
    \text{BW} = \sqrt{\frac{\sum_k (f_k - C)^2\, \bar{P}[k]}{P_{\text{tot}}}}
    \]

    \item \textbf{Entropia espectral} (\textit{specEntropy}): com $\tilde{P}[k] = \bar{P}[k] / P_{\text{tot}}$,
    \[
    H = -\sum_k \tilde{P}[k] \log_2\!\left(\tilde{P}[k] + \varepsilon\right)
    \]
    onde $\varepsilon$ é uma constante de estabilidade numérica.
\end{itemize}

\subsubsection{Características por Região Temporal}

Para cada uma das três regiões temporais, são calculados o centroide espectral, a entropia
espectral e a energia, de forma análoga às características globais acima descritas. Isto resulta
em nove características adicionais:
\textit{specCentroid\_start/mid/end},
\textit{specEntropy\_start/mid/end} e
\textit{energySpec\_start/mid/end}.

\subsubsection{Visualização e Seleção das Melhores Características}

As figuras seguintes mostram a distribuição de cada característica global por dígito.

% -----------------------------------------------------------------------
% FIGURA 2: Energia espectral por dígito (scatter plot)
% Gerada pelo bloco print_stft - figura "Distribuição da Energia Espectral Média por Dígito"
% Nome sugerido: stft_energia.png
% -----------------------------------------------------------------------
\begin{figure}[H]
    \centering
    \includegraphics[width=0.8\textwidth]{stft_energia.png}
    \caption{Distribuição da energia espectral por dígito.}
    \label{fig:stft_energia}
\end{figure}

% -----------------------------------------------------------------------
% FIGURA 3: Pico de frequência por dígito (scatter plot)
% Gerada pelo bloco print_stft - figura "Pico de Frequência por Dígito"
% Nome sugerido: stft_pico_freq.png
% -----------------------------------------------------------------------
\begin{figure}[H]
    \centering
    \includegraphics[width=0.8\textwidth]{stft_pico_freq.png}
    \caption{Distribuição do pico de frequência por dígito.}
    \label{fig:stft_pico_freq}
\end{figure}

% -----------------------------------------------------------------------
% FIGURA 4: Entropia espectral por dígito (scatter plot)
% Gerada pelo bloco print_stft - figura "Entropia Espectral"
% Nome sugerido: stft_entropia.png
% -----------------------------------------------------------------------
\begin{figure}[H]
    \centering
    \includegraphics[width=0.8\textwidth]{stft_entropia.png}
    \caption{Distribuição da entropia espectral por dígito.}
    \label{fig:stft_entropia}
\end{figure}

% -----------------------------------------------------------------------
% FIGURA 5: Centroide espectral por dígito (scatter plot)
% Gerada pelo bloco print_stft - figura "Centroide Espectral"
% Nome sugerido: stft_centroide.png
% -----------------------------------------------------------------------
\begin{figure}[H]
    \centering
    \includegraphics[width=0.8\textwidth]{stft_centroide.png}
    \caption{Distribuição do centroide espectral por dígito.}
    \label{fig:stft_centroide}
\end{figure}

% -----------------------------------------------------------------------
% FIGURA 6: Largura de banda espectral por dígito (scatter plot)
% Gerada pelo bloco print_stft - figura "Largura de Banda Espectral"
% Nome sugerido: stft_bw.png
% -----------------------------------------------------------------------
\begin{figure}[H]
    \centering
    \includegraphics[width=0.8\textwidth]{stft_bw.png}
    \caption{Distribuição da largura de banda espectral por dígito.}
    \label{fig:stft_bw}
\end{figure}

Com base na análise visual dos gráficos, as três características tempo-frequência que apresentam
maior capacidade de discriminação entre dígitos são o \textbf{centroide espectral}, a
\textbf{entropia espectral} e a \textbf{energia espectral}:

\begin{itemize}
    \item O \textbf{centroide espectral} reflete onde se concentra a maior parte da energia em
    frequência. Dígitos com sons vogais fechados (como ``one'' ou ``nine'') tendem a apresentar
    centroides em frequências distintas dos dígitos com vogais abertas (como ``zero'' ou ``four''),
    tornando esta característica informativa.

    \item A \textbf{entropia espectral} mede a irregularidade da distribuição espectral.
    Sons com espectros mais concentrados (tons quase periódicos) apresentam baixa entropia,
    enquanto sons com espectros planos ou ruidosos apresentam entropia elevada. Esta variação
    permite distinguir dígitos com conteúdo espectral mais ou menos estruturado.

    \item A \textbf{energia espectral} total do espectrograma complementa as anteriores ao
    capturar a intensidade global do sinal, sendo especialmente útil para separar dígitos
    com diferentes durações efetivas após o pré-processamento.
\end{itemize}

Todas as características foram guardadas na estrutura de dados para utilização posterior no
classificador.

\subsection{Transformada de Wavelet Discreta (DWT)}

A Transformada de Wavelet Discreta (DWT) foi aplicada como alternativa à STFT para a análise
tempo-frequência, permitindo uma resolução adaptativa: melhor resolução temporal para
altas frequências e melhor resolução em frequência para baixas frequências.

A wavelet utilizada foi a Daubechies 4 (\texttt{db4}), com um nível de decomposição $L=1$.
Para cada sinal $x[n]$, a DWT de nível 1 produz:

\begin{itemize}
    \item Coeficientes de \textbf{aproximação} $cA$: resultado da filtragem passa-baixo seguida
    de subamostragem, capturando as componentes de baixa frequência do sinal;
    \item Coeficientes de \textbf{detalhe} $cD$: resultado da filtragem passa-alto seguida de
    subamostragem, capturando as componentes de alta frequência.
\end{itemize}

A partir destes coeficientes, foram calculadas as energias:

\[
E_{\text{aprox}} = \sum_n cA[n]^2, \qquad E_{\text{detalhe}} = \sum_n cD[n]^2.
\]

Estas duas grandezas foram guardadas na estrutura de dados como \texttt{AproxEnergy} e
\texttt{DetailEnergy}.

% -----------------------------------------------------------------------
% FIGURA 7: Energia de aproximação e de detalhe DWT por dígito (subplot 1x2)
% Gerada pelo bloco %% 18 do main2_main.m
% Subplot esquerdo: energia de aproximação (pontos azuis)
% Subplot direito: energia de detalhe (pontos vermelhos)
% Nome sugerido: dwt_energias.png
% -----------------------------------------------------------------------
\begin{figure}[H]
    \centering
    \includegraphics[width=\textwidth]{dwt_energias.png}
    \caption{Energia dos coeficientes de aproximação (esquerda) e de detalhe (direita) obtidos
             pela DWT (wavelet \texttt{db4}, nível 1) para cada dígito.}
    \label{fig:dwt}
\end{figure}

\subsubsection{Comparação entre STFT e DWT}

Comparando os resultados obtidos com a STFT e com a DWT, verifica-se que:

\begin{itemize}
    \item A \textbf{STFT} oferece uma visão detalhada da evolução temporal do espectro, com
    resolução uniforme em frequência. O espectrograma é intuitivo e permite identificar visualmente
    padrões de cada dígito. As características extraídas (centroide, entropia, energia por região)
    apresentam boa discriminação.

    \item A \textbf{DWT} com a \texttt{db4} é mais adequada para capturar componentes transitórias
    do sinal (como plosivas e fricativas). A energia de detalhe tende a ser mais sensível a eventos
    rápidos, enquanto a energia de aproximação caracteriza melhor o conteúdo harmónico de baixa
    frequência.

    \item Em termos de discriminação por dígito, tanto a energia de aproximação como a de detalhe
    apresentam distribuições com alguma sobreposição entre classes, o que sugere que a DWT sozinha
    não é suficiente para uma classificação precisa, mas constitui um complemento útil às
    características da STFT.
\end{itemize}

\subsection{Classificação Automática}

\subsubsection{Seleção de Características}

Para a classificação, foram utilizadas as características que, combinadas, apresentaram melhor
capacidade de discriminação com base na análise visual efetuada nos pontos anteriores.

O conjunto de características final utilizado no classificador foi:

\begin{itemize}
    \item \textbf{Entropia espectral} (\textit{specEntropy}) -- característica tempo-frequência;
    \item \textbf{Largura de banda espectral} (\textit{specBW}) -- característica tempo-frequência;
    \item \textbf{Energia temporal na divisão 1} (\textit{energy\_int1}) -- característica temporal
    da Meta 1;
    \item \textbf{Energia temporal na divisão 2} (\textit{energy\_int2}) -- característica temporal
    da Meta 1.
\end{itemize}

Estas características foram normalizadas com \textit{z-score} antes de serem passadas ao
classificador:

\[
z = \frac{x - \mu}{\sigma}
\]

onde $\mu$ e $\sigma$ são, respetivamente, a média e o desvio padrão calculados sobre o conjunto
de treino.

\subsubsection{Modelo de Classificação: KNN}

O classificador utilizado foi o \textbf{K-Nearest Neighbors (KNN)} com $k=4$ vizinhos e distância
euclidiana, implementado com a função \texttt{fitcknn} do MATLAB. O conjunto de dados foi dividido
em 70\% para treino e 30\% para teste, com partição estratificada (\texttt{HoldOut}).

O KNN classifica cada amostra de teste com base nos $k$ exemplos de treino mais próximos no espaço
das características. A distância euclidiana entre duas amostras $\mathbf{x}$ e $\mathbf{y}$ é:

\[
d(\mathbf{x}, \mathbf{y}) = \sqrt{\sum_{i=1}^{n}(x_i - y_i)^2}.
\]

\subsubsection{Resultados da Classificação}

% -----------------------------------------------------------------------
% FIGURA 8: Matriz de confusão (confusionchart)
% Gerada pelo bloco %% 20 do main2_main.m
% Nome sugerido: knn_confusion_chart.png
% -----------------------------------------------------------------------
\begin{figure}[H]
    \centering
    \includegraphics[width=0.8\textwidth]{knn_confusion_chart.png}
    \caption{Matriz de confusão do classificador KNN ($k=4$), com indicação do número de amostras
             por célula.}
    \label{fig:confusion_chart}
\end{figure}


\subsubsection{Análise de Resultados}

A taxa de acerto global obtida pelo classificador KNN ($k=4$) no conjunto de teste (30\% dos dados)
foi de \textbf{74.00\%}, correspondendo a 111 amostras corretamente classificadas em 150.

\[
\text{Accuracy} = \frac{111}{150} \times 100 = 74{,}00\%
\]

Analisando os resultados por dígito, verificam-se diferenças consideráveis de desempenho entre
classes. O dígito \textbf{6} (\textit{six}) foi o melhor classificado, com uma taxa de acerto de
93,3\% (14 em 15 amostras), seguido dos dígitos \textbf{2} (\textit{two}) e \textbf{9}
(\textit{nine}), ambos com 86,7\%. Por outro lado, o dígito \textbf{0} (\textit{zero}) revelou-se
o mais difícil de classificar, com apenas 46,7\% de acerto (7 em 15), seguido do dígito \textbf{3}
(\textit{three}) com 53,3\%.

A confusão mais frequente ocorre entre os dígitos \textbf{3} e \textbf{0}: o dígito 3 foi
classificado como 0 em 6 das 15 amostras, e o dígito 0 foi classificado como 3 em 3 ocasiões.
Esta confusão bidirecional sugere que as características extraídas não conseguem separar
eficazmente estes dois dígitos, provavelmente por partilharem padrões espectrais e de energia
semelhantes. Outras confusões notáveis incluem o dígito \textbf{5} classificado como 7 (3 vezes)
e o dígito \textbf{8} classificado como 2 (3 vezes), o que é consistente com a semelhança
fonética entre estes pares em inglês.

\section{Conclusões}

Nesta Meta 2 foi desenvolvida a componente de análise tempo-frequência e classificação do
\textit{pipeline} de reconhecimento de dígitos falados.

A \textbf{STFT} revelou-se uma ferramenta eficaz para visualizar e caracterizar a evolução
espectral dos sinais. A escolha de parâmetros (janela de 256 amostras, sobreposição de 50\%,
NFFT=256) foi validada empiricamente e produz um bom equilíbrio entre resolução temporal e
espectral. As características extraídas -- centroide, entropia e energia espectral -- mostraram
capacidade discriminativa razoável entre classes.

A \textbf{DWT} com a wavelet \texttt{db4} complementou a análise ao capturar informação de
detalhe de alta frequência. Apesar de a energia de detalhe apresentar alguma separação entre
dígitos, esta técnica sozinha não supera a STFT na discriminação. A combinação de ambas pode
ser explorada em trabalho futuro.

O \textbf{classificador KNN} com $k=4$ e as quatro características selecionadas atingiu uma taxa
de acerto que demonstra a viabilidade da abordagem. Os erros de classificação mais frequentes
ocorrem entre dígitos com padrões acústicos semelhantes, o que é esperado dado que as
características utilizadas são de natureza estatística e não capturam toda a informação fonética
do sinal.

Como propostas de melhoria, destacam-se:
\begin{itemize}
    \item Utilização de características de ordem superior (ex.\ coeficientes MFCC -- Mel-Frequency
    Cepstral Coefficients), que são especificamente desenhados para representação de voz;
    \item Aplicação de modelos de classificação mais sofisticados, como \textit{Support Vector
    Machines} (SVM) ou redes neuronais;
    \item Aumento da decomposição DWT para múltiplos níveis, extraindo características de cada
    sub-banda;
    \item Uso de validação cruzada mais robusta (ex.\ \textit{k-fold} com $k$ maior) para uma
    estimativa mais estável do desempenho.
\end{itemize}

No geral, o trabalho desenvolvido ao longo das duas metas permitiu construir um sistema completo
de análise e classificação de sinais de voz, desde a aquisição e pré-processamento até à
classificação automática, explorando ferramentas fundamentais da análise de sinais no tempo,
frequência e tempo-frequência.

\vspace{2em}
\noindent\rule{\textwidth}{0.4pt}
\vspace{0.5em}

\noindent\small\textit{Nota: Este relatório foi redigido pelo autor com recurso a um modelo de
linguagem (LLM) para revisão, reformulação e reestruturação de algumas secções. Todo o código,
análise e resultados são da responsabilidade do autor.}

\end{document}
